% ID: SEL-D0033-OA-E02
% ESTADO: cotejada. Segunda pasada visual y borrador compilado revisados.
% Fuente: Selectividad/matematicas_aplicadas/titular_septiembre/Examen-Matemáticas Aplicadas a las Ciencias Sociales.pdf
% SHA-256: 0f95b6aefa78f1a13e2115b281338bd807177671c303c380180957fe03f0a605
% Página: 1; opción A; ejercicio original 2.

El consumo de cereales en una ciudad, en miles de toneladas, viene dado por la función $c(t)=t^3-15t^2+63t+10$, para $0\leq t\leq12$, donde $t$ representa el tiempo.
\begin{enumerate}[label=\alph*)]
\item \textbf{(0.8 puntos)} \textquestiondown{}En qué instante se alcanza el máximo consumo de cereales y cuántas toneladas se consumen en ese momento?
\item \textbf{(0.7 puntos)} \textquestiondown{}En qué intervalo de tiempo decrece el consumo de cereales?
\item \textbf{(1 punto)} Represente gráficamente la función.
\end{enumerate}
